\documentclass[a4paper,12pt]{scrartcl}
\author{Anisa Mecheraoui}

\usepackage[ngerman]{babel}
\usepackage[]{csquotes}
\usepackage{titling}
\usepackage{fancyhdr}
\usepackage{graphicx}
\usepackage{float}
\usepackage[margin=2.5cm,includeheadfoot, footskip=80pt]{geometry}
\usepackage[breaklinks]{hyperref}
\usepackage{blindtext}
\usepackage[nospace]{varioref}
\usepackage[backend=biber,style=apa6]{biblatex}
\usepackage{numprint}
\usepackage{pgfplots}
\pgfplotsset{compat=1.16}
\usepackage{booktabs}
\usepackage[justification=centering]{caption}
\usepackage[titles]{tocloft}
\usepackage{tabularx}
\usepackage{setspace}
\setstretch{1.5}
\usepackage{sectsty}
\sectionfont{\fontsize{14}{16.8}\selectfont}
\subsectionfont{\fontsize{12}{14.4}\selectfont}
\usepackage{amsmath}
\usepackage{pgffor}
\usepackage[roman]{parnotes}
\usepackage{xifthen}
\usepackage{units}
\usepackage{xcolor}

%bib referenc
\addbibresource{myBib.bib}
%graphicspath
\graphicspath{{images/}}
\definecolor{ba-blau}{HTML}{093a80}
\setlength{\columnsep}{1cm}
\newlength{\drop}
\newcommand{\myLarge}[1]{{\fontsize{20pt}{24pt}\selectfont\color{ba-blau}#1}}
\newcommand{\pageTitel}[1]{{\fontsize{14pt}{17pt}\selectfont\textbf{\color{ba-blau}#1}}}
\fancypagestyle{style1}{
    \setlength{\headheight}{10pt}
    \pagestyle{fancy}
    \fancyhead{}
    \setlength{\headheight}{2cm}
    \chead{\includegraphics[]{deckbild.jpeg} \\}
    %pagefooter
    \fancyfoot{}
    \fancyfoot[R]{\thepage}
}
\fancypagestyle{Deckblatt}{
    \setlength{\headheight}{10pt}
    \pagestyle{fancy}
    \fancyhead{}
    \setlength{\headheight}{2cm}
    \setlength{\headsep}{80pt}
    \chead{\includegraphics[]{deckbild.jpeg} \\}
    %pagefooter
    \fancyfoot{}
    \fancyfoot[R]{\thepage}
}
\addtokomafont{disposition}{\color{ba-blau}}
\addtokomafont{sectionentry}{\color{ba-blau}}

\begin{document}

\thispagestyle{Deckblatt}
%use different pagestyling
\pagenumbering{gobble}

%% title page
\begin{center}
    \myLarge {Titel der Bachelor Thesis} \\[3ex]
    \includegraphics[width=3cm]{Platzhalter_logo.png} \\[3ex]
    \large{Exposé erstellt im Rahmen des Bachelor Thesis Kolloquiums:
} \\
    \large{DD.MM.YYYY}
\end{center}

\renewcommand{\arraystretch}{2.5}
\begin{table}[h]
    \centering
    \begin{tabularx}{\textwidth}{l c}
        Studiengruppe: &  XX\_XX\_20XX-I \\
        Name Studierender: & Max Mustermann\\
        Anzahl der Wörter: & 0\\
        (inkl. wörtliche Zitate / Fußnoten) & ~ \\
        Anzahl der Wörter: & \\
        (exkl. wörtliche Zitate / Fußnoten) & ~ \\
        Akademischer Gutachter: & Prof. Mustermann\\
        Betrieblicher Gutachter: & Herr Firma\\
        Abgabedatum: & DD.MM.YYYY\\
    \end{tabularx}
\end{table}

\newpage

\pagestyle{style1}
\pageTitel{Gleichbehandlung der Geschlechter}\\
Aus Gründen der besseren Lesbarkeit wird bei Personenbezeichnungen und personenbezogenen Hauptwörtern in diesem Projektbericht die maskuline Form verwendet. Entsprechende Begriffe gelten im Sinne der Gleichbehandlung grundsätzlich für alle Geschlechter. Die verkürzte Sprachform hat nur redaktionelle Gründe und beinhaltet keine Wertung
\newpage
\pagenumbering{Roman}
\setcounter{page}{1}
\tableofcontents
\newpage
\listoffigures
\newpage
\listoftables
\newpage

%stores roman numbercounter
\newcounter{preamblecounterstate}
\setcounter{preamblecounterstate}{\value{page}}

% use arabic numbers for actual content pages
\pagenumbering{arabic}

%includes pages written in the folder section
%\section{Demo-Seite}
%Falls keine Seitenanzahl vorhanden ist  dann:\footcite[Vgl.][]{DemoQuelle}
Demo-Quelle \footcite[Vgl.][\printfield{pages}]{DemoQuelle}

\begin{center}
    \begin{table}[h]
    \centering
    \begin{tabular}{|c|p{6cm}|}
        \hline
        \textbf{Datum} & \textbf{Aktivitäten} \\
        \hline
        Kebab & 7€ \\
        \hline
        Adana & \begin{itemize}
            \item \textbf{Groß}: 8€
            \item \textbf{Klein}: 6€
        \end{itemize} \\
        \hline
        Köfte & \begin{itemize}
            \item 5 Stück: 8€
            \item 2 Stück: 6€
        \end{itemize}\\
        \hline
        Mercimek Suppe & 3€ \\
        \hline
        Dönerteller & 15€ \\
        \hline
    \end{tabular}
    \caption{Beispiel Tabelle.}
    \label{tab:example}
    \end{table}
    \end{center}

Die Tabelle zeigt den Preis eines Dönertellers, dieser lässt sich wie folgt berechnen:
\begin{equation}
    15 = \sum_{n=1}^{10} \frac{n}{20} + \sum_{k=1}^{5} \frac{2k}{10} - \sum_{i=1}^{3} i
\end{equation}
\newpage
Der folgende Abschnitt könnte hilfreich für eine Ausarbeitung in der Informatik sein.
\begin{lstlisting}
    def hello_world():
        print("Hello, World!")
\end{lstlisting}
\section{Kontext der Ausarbeitung}


\subsection{Relevanz der Untersuchung}


\subsection{Stand der Forschung}


\section{Erstellung der Arbeit}
Test\fullfootcite{goos_single_2002}
\subsection{Genutzte Quellen}
\subsection{Zeitplan}
\section{Zielsetzung}


\subsection{Grundlagen}


\subsection{Forschungsfrage}



\newpage

% Set page number to Roman for bibliography
\pagenumbering{Roman}

%reads counter stored earlier
\setcounter{page}{\value{preamblecounterstate}}

\printbibliography
\end{document}